% Tabulkové poznámky (sepfootnotes aparát tabfn) – definice obsahu

\tabfnnotecontent{vymezeni-aot40}{\textit{Pro účely tohoto zákona AOT40 znamená součet rozdílů mezi hodinovou koncentrací větší než 80 µg.m\textsuperscript{-3} (= 40 ppb) a hodnotou 80 µg.m\textsuperscript{-3} v dané periodě užitím pouze hodinových hodnot změřených každý den mezi 08:00 a 20:00 SEČ, vypočtený z hodinových hodnot v letním období (1. května - 31. července).} \parencite[příl.~1, pozn.~5]{wolterskluwerZakon20120122012}}

\tabfnnotecontent{vymezeni-la-a-index}{Index „A“ v označení $L_A$ znamená, že jde o hladinu akustického tlaku měřenou s frekvenční váhovou charakteristikou A (A-vážená hladina); např. $LA_{eq,T}$ je A-vážená ekvivalentní hladina akustického tlaku za dobu $T$. \parencite{info@wolterskluwer.czNarizeni2722011}}

\tabfnnotecontent{vymezeni-la-eu-note}{\textit{Poznámka k EU (proč tu nejsou hodnoty):} Směrnice \parencite{SmerniceEvropskehoParlamentu2002} nestanovuje jednotné číselné hygienické limity (ty si určují členské státy). Zavádí ale společné ukazatele hluku: $L_{den}$ (\textit{day-evening-night level}) je vážená průměrná hladina za 24 h, kde se večerní část penalizuje +5 dB a noční část +10 dB; $L_{night}$ (\textit{night level}) je průměrná hladina za noční období používaná hlavně pro hodnocení rušení spánku.
Zdroj: \parencite{info@wolterskluwer.czNarizeni2722011}.}
